%! AP Statistics — Unit 1, Lesson 1.4 — Exit Ticket ANSWER KEY
\documentclass[10pt]{article}
\usepackage{apstats-article}
\usepackage{apstats-key}

%=============================================================================
\begin{document}

\pageheader{Unit 1, Lesson 1.4}{Exit Ticket --- Answer Key}
\begin{tabularx}{\linewidth}{X X X}
Name:\;\ans{Key} & Date:\; & Period:\;
\end{tabularx}

\vspace{0.22in}

\begin{tcolorbox}[colback=sky, colframe=navy, boxrule=0.9pt, arc=2mm,
    left=4mm, right=4mm, top=3mm, bottom=3mm]
\small A survey of \textbf{80 adults}: preferred TV genre.
\end{tcolorbox}

\vspace{0.2in}

\renewcommand{\arraystretch}{2.2}
\begin{tabularx}{\linewidth}{@{} >{\bfseries}p{0.22\linewidth}
                                    C{0.18\linewidth}
                                    C{0.26\linewidth} @{}}
\rowcolor{sky}
\textbf{Genre} & \textbf{Frequency} & \textbf{Rel.\ Frequency} \\
\midrule
News    & 32 & \ans{0.40 (40\%)} \\
Sports  & 20 & \ans{0.25 (25\%)} \\
Reality & 16 & \ans{0.20 (20\%)} \\
Other   & 12 & \ans{0.15 (15\%)} \\
\midrule
\rowcolor{sky}\textbf{Total} & \ans{80} & \ans{1.00} \\
\end{tabularx}

\vspace{0.24in}

\begin{enumerate}[leftmargin=1.5em, itemsep=10pt]

  \item \ans{Of the 80 adults surveyed, news was the most preferred TV genre
        (40\%), while ``other'' was the least preferred (15\%). Sports (25\%)
        and reality TV (20\%) fell in between.}

  \begin{teachernote}
  Require: $n = 80$, variable in context, most and least common category with
  relative frequencies. Deduct for missing context or for giving only counts.
  \end{teachernote}

  \item \ans{If the y-axis started at 10\% instead of 0, the News bar would
        appear roughly three times taller than the Other bar, even though the
        actual ratio is $40\%/15\% \approx 2.7\times$. A reader might incorrectly
        conclude that news is much more dominant than the data actually show.}

  \begin{teachernote}
  Accept any of the three tricks. The response must (1) name the specific
  misleading feature and (2) state the incorrect conclusion a reader would draw.
  \end{teachernote}

\end{enumerate}

\vspace{0.18in}

\begin{tcolorbox}[colback=redbg, colframe=redacc, boxrule=0.8pt, arc=2mm,
    left=4mm, right=4mm, top=3mm, bottom=3mm]
\small\textbf{AP-Style Practice --- Key}

\vspace{8pt}
\ans{The student's statement is not fully correct. While both charts use the
same categories (making the labels comparable), using different scales on the
y-axis makes a direct visual comparison invalid. School A's bars show counts
(which depend on total enrollment), while School B's bars show relative
frequencies (proportions). A student at a larger school is not necessarily
represented by a taller bar. To compare the two schools fairly, both charts
should use relative frequency on the y-axis, or the data should be presented
in the same type of display.}

\begin{teachernote}
AP rubric would award credit for: (1) correctly identifying the incompatibility
of the two y-axes, (2) explaining why counts and relative frequencies cannot be
compared directly, (3) suggesting a fix. A response that only says ``they use
different y-axes'' without explaining the implication earns partial credit.
\end{teachernote}
\end{tcolorbox}

\end{document}
